\renewcommand{\Anno}{1998}

\renewcommand{\NumeroProblema}{5}

\renewcommand{\Testo}{Dato un numero dispari $n$, si consideri il
seguente algoritmo: si calcoli
$$a_{1}=\frac{3n+1}{2};$$
se $a_{1}$ è pari l'algoritmo si arresta, altrimenti si calcoli
$$a_{2}=\frac{3a_{1}+1}{2};$$
se $a_{2}$ è pari l'algoritmo si arresta, altrimenti si calcoli
$$a_{3}=\frac{3a_{2}+1}{2};$$
e cos\`{\i} via.
	
Dimostare che, qualunque sia il numero $n$ considerato, l'algoritmo prima o poi si arresta, cioè la successione $a_{1},a_{2},\ldots$ da esso generata è finita.
	
Determinare inoltre da quanti termini essa è costituita.}

\renewcommand{\Risposta}{Scritto $n$ nella forma $2^{k}d-1$, con $d$
dispari, la successione risulta composta da $k$ termini.}

\renewcommand{\Aiutino}{Si consiglia di familiarizzare con il problema
eseguendo l'algoritmo a partire da valori bassi di $n$ (provarli tutti fino a 20), segnandosi il numero di passi necessario, e quindi
formulare una congettura da dimostrare.  }

\renewcommand{\SoluzioneUfficiale}{
}

\renewcommand{\SoluzioniAlternative}{
}

\renewcommand{\Osservazioni}{
}

\renewcommand{\Commenti}{
}
% Aggiunta redazionale: la soluzione preesistente non viene modificata.
\providecommand{\SoluzioneRedazionale}{}
\renewcommand{\SoluzioneRedazionale}{
% SSSUP-BOZZA-AI: SoluzioneRedazionale; autore=ChatGPT; data=2026-11-08; stato=DA-VERIFICARE
% SSSUP-PROVENIENZA: AI: elaborazione nuova dal testo
Per un intero dispari $n\neq-1$, scriviamo
\[
n+1=2^k d,\qquad k\geq1,\quad d\text{ dispari}.
\]
Posto $a_0=n$, la relazione dell'algoritmo equivale a
\[
a_{j+1}+1=\frac32(a_j+1).
\]
Per induzione, finché l'algoritmo prosegue,
\[
a_j=3^j2^{k-j}d-1\qquad(0\leq j\leq k).
\]
Per $j<k$ il termine a destra è dispari; per $j=k$ è pari. Pertanto vengono prodotti esattamente $k$ termini $a_1,\ldots,a_k$.

\emph{Precisazione sul dominio.} La conclusione vale in particolare per ogni intero positivo dispari. Se ``numero dispari'' comprende anche tutti gli interi negativi, l'enunciato ha l'eccezione $n=-1$: in quel caso tutti i termini sono $-1$ e l'algoritmo non si arresta. Il testo originale non è stato modificato.
}
